% French PDF pp. 31--36; proof of 4.3.10 continues from en-05.
Now consider \(f=i_Rg\) and the fibered product \(R'=R\times_S T\):
\nde{Denote by \(p\) the projection \(R'\to R\). Since \(i_R\) is a
monomorphism, one has \(p=gi_{R'}\). Let \(g'\) be the section of
\(i_{R'}\) defined by \(g\); then \(pg'=g\), hence
\(pg'i_{R'}=gi_{R'}=p\); this implies \(g'i_{R'}=\mathrm{id}_{R'}\),
and thus \(i_{R'}:R'\to T\) is an isomorphism, with inverse \(g'\).}
\[
\xymatrix@R=2.5em{
P\ar[r]^{\ell_P} & LP\\
R\ar[u]^u\ar[r]^{i_R} & S\ar[u]_{z_R(u)}\\
R'\ar[u]^p\ar[r]_{i_{R'}}^{\sim}
 & T\ar[u]_f\ar[ul]_g }.
\]
By definition of \(\ell_P\), \(\ell_Pug=z_R(u)f\) (this is the
particular case of what one is trying to prove in which \(i_R\) is
an isomorphism), and \(z_R(u)f=z_R(u)i_Rg\).

(ii) and (iv) merely express the definition of \(\Hom(S,LP)\) as an
inductive limit.

(iii): If \(K\) denotes the kernel of the pair \((u,v)\), for each
morphism \(f:S\to Q\), where \(S\) is representable, \(K\times_Q S\)
is a subfunctor of the kernel of the pair of arrows \(uf,vf:S\to P\).
By (T' 0), one is therefore reduced to proving the assertion in the
case where \(Q=S\) is representable. But in this case it follows from
(ii) and (iv) that \(K\) contains a refinement of \(S\), hence is a
refinement of \(S\).
\end{proof}

Finally one verifies that \(P\mapsto LP\) defines a functor
\[
L:\widehat{\mathcal C}\longrightarrow\widehat{\mathcal C}
\]
and \(P\mapsto\ell_P\) a morphism of functors
\[
\ell:\mathrm{Id}_{\widehat{\mathcal C}}\longrightarrow L.
\]
Let us now state the essential result:

\begin{proposition}{4.3.11}\label{IV.4.3.11}
(i) If \(P\) is any presheaf, \(LP\) is separated and
\(\ell_P:P\to LP\) is covering (4.2.3).

(ii) If \(P\) is a sheaf, \(P\to LP\) is an isomorphism.

(iii) For every presheaf \(P\) and every separated presheaf
(resp. sheaf) \(F\), the map
\[
\Hom(\ell_P,F):\Hom(LP,F)\longrightarrow\Hom(P,F)
\]
is injective (resp. bijective).

(iv) If \(P\) is separated, \(\ell_P:P\to LP\) is a covering
monomorphism (hence \(P\) is a refinement of \(LP\)), and \(LP\)
is a sheaf.
\end{proposition}
\begin{proof}
(i) First, \(P\to LP\) is covering; indeed, for every morphism
\(S\xrightarrow{v}LP\), there exists, by 4.3.10 (i) and (ii), a
refinement \(R\) of \(S\) such that one has a commutative diagram:
\[
\xymatrix@R=2.3em{
 & P\ar[r]^{\ell_P} & LP\\
 & P\times_{LP}S\ar[u]^{v'}\ar[r] & S\ar[u]_v\\
R\ar[uur]\ar[ur]\ar[urr]_{i_R} & & },
\]
hence \(P\times_{LP}S\to S\) is covering. It follows, by (C 0),
that \(P\to LP\) is covering.

If, on the other hand, two morphisms \(S\xrightarrow{v_1,v_2}LP\)
induce the same morphism from a refinement \(R\) of \(S\), let us show
that they are equal. There exist refinements \(R_i\), \(i=1,2\), and
morphisms \(u_i:R_i\to P\) such that \(z_{R_i}(u_i)=v_i\).
Taking \(R\) sufficiently small, one may assume that \(R_1=R_2=R\).
It then follows from the commutative diagram of 4.3.10 (i) that
\(\ell_Pu_1=\ell_Pu_2\). By \loccit{} (iii), \(u_1\) and \(u_2\)
therefore agree on a refinement of \(R\), hence a refinement of \(S\),
which implies \(z_R(u_1)=z_R(u_2)\), by \loccit{} (iv).

(ii) is clear, for if \(P\) is a sheaf,
\(\Hom(S,P)\to\Hom(R,P)\) is already an isomorphism for every
refinement \(R\) of \(S\).

(iii) Let \(u,v\) be two morphisms \(LP\to F\) such that
\(u\ell_P=v\ell_P\). To show that \(u=v\), it suffices to see that
\(uf=vf\) for every \(f:S\to LP\), where \(S\) is representable.
Now there exists a refinement \(R\) of \(S\) and a morphism
\(g:R\to P\) with \(f=z_R(g)\). Then \(uf\) and \(vf\) agree on \(R\)
with \(u\ell_Pg=v\ell_Pg\), hence agree on \(R\).
If \(F\) is separated, one therefore has \(uf=vf\).
Suppose now that \(F\) is a sheaf; one then has the commutative diagram
\[
\xymatrix{
P\ar[r]^{\ell_p}\ar[d] & LP\ar[d]\ar[dl]\\
F\ar[r]_{\sim} & LF }
\]
which shows that \(\Hom(\ell_P,F)\) is surjective.

(iv) Let us show that if \(P\) is separated, \(P\to LP\) is a
monomorphism. For this, it suffices to see that for every pair of
morphisms \(u,v:S\to P\) (where \(S\) is representable) such that
\(\ell_Pu=\ell_Pv\), one has \(u=v\).
Now \loccit{} (iii) shows that \(u\) and \(v\) agree on a refinement
of \(S\), hence agree, since \(P\) is separated.
This shows that \(P\to LP\) is a monomorphism; since it is covering
by (i), one obtains that \(P\) is a refinement of \(LP\).

Finally let us show that \(LP\) is a sheaf.
Since one already knows by (i) that it is a separated presheaf,
it suffices to see that for every \(S\in\Ob\mathcal C\), every
refinement \(R\) of \(S\), and every morphism \(h:R\to LP\), there
exists a morphism \(u:S\to LP\) with \(ui_R=h\).
Now \(R'=P\times_{LP}R\) is a refinement of \(R\), since \(P\) is
a refinement of \(LP\), hence \(R'\) is a refinement of \(S\).
Set \(u=z_{R'}(h')\):
\[
\xymatrix{
P\ar[r]^{\ell_p} & LP\\
R'\ar[u]^{h'}\ar[r]_j & R\ar[u]^h\ar[r]_{i_R} & S\ar[ul]^u }.
\]
One has \(ui_{R'}=\ell_Ph'=hj\), hence \(ui_Rj=hj\).
Since \(R'\) is a refinement of \(R\) and \(LP\) is separated,
4.3.2 shows that \(ui_R=h\).
\end{proof}

\begin{corollary}{4.3.12}\label{IV.4.3.12}
\nde{The statement of the corollary and its proof have been given in detail.}
Let \(F\) be a sheaf and \(R\) a sub-\(\widehat{\mathcal C}\)-object
of \(F\). Then \(R\) is a separated presheaf,
\(\ell_R:R\to LR\) is a covering monomorphism, and one has a
commutative diagram
\[
\xymatrix{
R\ar@{^{(}->}[rr]^i\ar@{^{(}->}[dr]_{\ell_R} & & F\\
 & LR\ar@{^{(}->}[ur]_j & }.
\]
Consequently, \(R\) is a refinement of \(F\) if and only if \(j\)
is an isomorphism.
\end{corollary}
\begin{proof}
One has already noted that \(R\) is separated and that \(j\) is a
monomorphism, cf. N.D.E. (24) and (26).
By 4.3.11 (iv), \(\ell_R\) is a covering monomorphism.
Thus, if \(j\) is an isomorphism, \(R\) is a refinement of \(F\).
Conversely, if \(i\) is covering, \(j\) is too, hence is an
isomorphism by 4.3.2.1.
\end{proof}

\begin{remark}{4.3.13}\label{IV.4.3.13}
If \(J'(S)\) is a cofinal subset of \(J(S)\), one has
\[
\Hom(S,LP)=\varinjlim_{R\in J'(S)}\Hom(R,P).
\]
In particular, let \(S\mapsto R(S)\) be a pretopology generating
the given topology. The functor \(L\) may be described with the aid
of the covering families that are elements of \(R(S)\).
By making the formula above explicit, one recovers the construction of [MA].
\end{remark}

Denote by \(i\) the inclusion functor
\(\widetilde{\mathcal C}\to\widehat{\mathcal C}\).
The following theorem follows from Proposition 4.3.11:

\begin{theorem}{4.3.14}\label{IV.4.3.14}
There exists a unique functor \(a:\widehat{\mathcal C}\to
\widetilde{\mathcal C}\) such that the following diagram is commutative
\[
\xymatrix{
\widehat{\mathcal C}\ar[r]^L\ar[d]_a & \widehat{\mathcal C}\ar[d]^L\\
\widetilde{\mathcal C}\ar[r]_i & \widehat{\mathcal C} },
\]
i.e. for every presheaf \(P\), \(L(L(P))\) is a sheaf.
The functors \(i,a\) are adjoint to one another: for every presheaf
\(P\) and every sheaf \(F\), one has an isomorphism functorial in \(P,F\)
\[
\Hom_{\widehat{\mathcal C}}(P,i(F))
\simeq\Hom_{\widetilde{\mathcal C}}(a(P),F),
\]
that is,
\[
\Hom(P,F)\simeq\Hom(a(P),F).
\]
\end{theorem}

\begin{definition}{4.3.15}\label{IV.4.3.15}
The sheaf \(a(P)\) is said to be \emph{associated} to the presheaf \(P\).
\end{definition}
\begin{remark}{4.3.16}\label{IV.4.3.16}
Since the functor \(L\) is constructed with the aid of projective
limits and filtered inductive limits, it commutes with finite
projective limits.\nde{In particular, if \(K\) is the kernel of a
pair of morphisms of presheaves \(u,v:Q\rightrightarrows P\),
then \(LK\) is the kernel of \(Lu,Lv:LQ\rightrightarrows LP\)
(this will be used in 4.4.5).}
Moreover, if one identifies \(L(P\times P)\) with \(LP\times LP\),
the morphism \(\ell_{P\times P}\) is identified with
\(\ell_P\times\ell_P\). It follows, for example, that if \(P\) is a
presheaf of groups, \(LP\) is also canonically endowed with the
structure of a presheaf of groups and the canonical morphism
\(P\to LP\) is a morphism of groups. The same holds for the functor
\(a\), which shows that if \(P\) is a presheaf of groups and \(F\)
a sheaf of groups, one has an isomorphism
\[
\Hom_{\widehat{\mathcal C}\text{-gr.}}(P,i(F))
\simeq\Hom_{\widetilde{\mathcal C}\text{-gr.}}(a(P),F).
\]
See [D] for more details.
\end{remark}

\numpara{4.3.17}\label{IV.4.3.17}
If \(\mathcal V\) is any category, one calls a \emph{presheaf on
\(\mathcal C\) with values in \(\mathcal V\)} a contravariant
functor from \(\mathcal C\) to \(\mathcal V\).
To define sheaves with values in \(\mathcal V\), we must first
recall the definition of the \emph{projective limit} of a functor.
If \(\mathcal R,\mathcal V\) are two categories, and
\[
F:\mathcal R^\circ\longrightarrow\mathcal V
\]
is a contravariant functor from \(\mathcal R\) to \(\mathcal V\),
one denotes by \(\varprojlim F\) the object of
\(\widehat{\mathcal V}\) defined as follows:
% typo?: kept as printed: Hom_V(F(U),X) in the displayed formula.
\[
(\varprojlim F)(X)=\Hom_{\widehat{\mathcal V}}(X,\varprojlim F)
=\varprojlim_{U\in\Ob\mathcal R}\Hom_{\mathcal V}(F(U),X)
=\Hom(c_X,F),
\]
where \(X\) is a variable object of \(\mathcal V\), where \(c_X\)
denotes the contravariant functor from \(\mathcal R\) to
\(\mathcal V\) which sends each object of \(\mathcal R\) to \(X\)
and each arrow of \(\mathcal R\) to \(\mathrm{id}_X\), and where
the last \(\Hom\) is taken in the category
\(\Hom(\mathcal R^\circ,\mathcal V)\).
If \(\mathcal R\) possesses a final object \(e_{\mathcal R}\),
one has \(\varprojlim F=F(e_{\mathcal R})\).
If \(\mathcal V\) is the category of sets, the functor
\(\varprojlim\) is identified with the functor \(\Gamma\).

If \(S\) is an object of \(\mathcal C\) and \(R\) a sieve of \(S\),
denote by \(\mathcal R\) the full subcategory of \(\mathcal C_{/S}\)
whose set of objects is \(E(R)\), and by
\(i_R:\mathcal R\to\mathcal C_{/S}\) the canonical functor.
If \(P\) is a presheaf on \(\mathcal C\) with values in
\(\mathcal V\), it defines by restriction a functor
\(P_S:(\mathcal C_{/S})^\circ\to\mathcal V\).
The functor \(i_R\) induces a morphism of \(\widehat{\mathcal V}\):
\[
P(S)=P_S(S)=\varprojlim P_S\longrightarrow\varprojlim(P_S\circ i_R).
\]
One denotes by \(\varprojlim_R P\) the object
\(\varprojlim(P_S\circ i_R)\) of \(\widehat{\mathcal V}\).
By 4.1.6, Definition 4.3.1 generalizes to the

\begin{definition}{4.3.18}\label{IV.4.3.18}
The presheaf \(P\) on \(\mathcal C\) with values in \(\mathcal V\)
is said to be \emph{separated} (resp. a \emph{sheaf}) if, for
every \(S\in\Ob\mathcal C\) and every \(R\in J(S)\), the
canonical morphism of \(\widehat{\mathcal V}\)
\[
P(S)\longrightarrow\varprojlim_R P
\]
is a monomorphism (resp. an isomorphism).
\end{definition}

In the case where \(\mathcal V\) is the category (Gr.) of groups
(or any other category of sets endowed with algebraic structures
defined by finite projective limits), one can see (cf. [D]) that
there is an equivalence between the following notions: a presheaf
on \(\mathcal C\) with values in (Gr.) whose underlying presheaf
of sets is a sheaf, and a group in the category of sheaves of sets.
Taking these identifications into account, we shall always
consider sheaves with values in a category of sets endowed with
algebraic structures defined by finite projective limits as
sheaves of sets endowed, in the category \(\widetilde{\mathcal C}\),
with the corresponding algebraic structure.

\sgasection{4.4}{Exactness properties of the category of sheaves}
\begin{theorem}{4.4.1}\label{IV.4.4.1}
(i) Arbitrary projective limits exist in \(\widetilde{\mathcal C}\);
\og they are calculated in \(\widehat{\mathcal C}\)\fg, i.e. the
inclusion functor \(i:\widetilde{\mathcal C}\to\widehat{\mathcal C}\)
commutes with projective limits: if \((X_\alpha)\) is a projective
system of sheaves, the presheaf
\[
\varprojlim i(X_\alpha):S\longmapsto\varprojlim X_\alpha(S)
\]
is a sheaf, and one has
\(i(\varprojlim X_\alpha)=\varprojlim i(X_\alpha)\).

(ii) Arbitrary inductive limits exist in \(\widetilde{\mathcal C}\):
if \((X_\alpha)\) is an inductive system of sheaves, one has
\[
\varinjlim X_\alpha=a(\varinjlim i(X_\alpha)),
\]
where \(\varinjlim i(X_\alpha)\) is the inductive-limit presheaf of
the \(i(X_\alpha)\):
\[
\varinjlim i(X_\alpha):S\longmapsto\varinjlim X_\alpha(S).
\]

(iii) The functor \(a:\widehat{\mathcal C}\to\widetilde{\mathcal C}\)
commutes with arbitrary inductive limits and finite projective limits.
\end{theorem}
\begin{proof}
Assertions (i) and (ii) follow formally from the adjunction formula
(4.3.14), and\nde{The original has been modified here.}
the first assertion of (iii) follows from (ii).
Finally, the second assertion of (iii) has already been pointed
out in 4.3.16.
\end{proof}

\begin{scholium}{4.4.2}\label{IV.4.4.2}
This theorem allows one to use the following method to prove in
\(\widetilde{\mathcal C}\) an assertion concerning simultaneously
arbitrary inductive limits and finite projective limits (for example:
\og every epimorphism is universal effective\fg, cf. below).
One begins by proving the corresponding assertion in the category
of sets, then extends it \og argument by argument\fg{} to the
category of presheaves. Next, one uses the preceding theorem to
pass from the category of presheaves to the category of sheaves.
One will see many examples of this method in what follows
(4.4.3, 4.4.6, 4.4.9, etc.).
\end{scholium}

Finally let us remark that the assertions concerning the category
of presheaves are formally corollaries of the assertions concerning
the category of sheaves. Indeed, it suffices to take as topology
the least fine topology (\og chaotic\fg), that is, the topology
defined by \(J(S)=\{S\}\) for every \(S\in\Ob\mathcal C\);
every functor is indeed a sheaf for this topology.

\begin{proposition}{4.4.3}\label{IV.4.4.3}
Let \(\mathcal F=\{F_i\to F\}\) be a family of morphisms of sheaves.
The following conditions are equivalent:
\begin{enumeratei}
\item \(\mathcal F\) is an epimorphic family.
\item \(\mathcal F\) is a universal effective epimorphic family (1.13).
\item \(\mathcal F\) is covering (4.2.3).
\item The image sheaf of \(\mathcal F\) (that is, the sheaf associated
to the image presheaf of \(\mathcal F\) (4.1.2)) is \(F\).
\end{enumeratei}
\end{proposition}
The equivalence of (iii) and (iv) follows from 4.3.12.
The other equivalences will follow from the following lemmas.

\begin{lemma}{4.4.4}\label{IV.4.4.4}
Let \(f:X\to Y\) be a monomorphism of sheaves which is an
epimorphism. Then \(f\) is an isomorphism.
\end{lemma}
\begin{proof}
The lemma is first clear in the category of sets.
Let us next prove it in the category of presheaves.
Consider the presheaf
\[
V:S\longmapsto Y(S)\amalg^{X(S)}Y(S);
\]
it is the amalgamated sum of \(Y\) and \(Y\) under \(X\) in the
category of presheaves.\nde{The rest of the proof has been
slightly modified.} Denote by \(i_1,i_2\) the two
\og coordinate\fg{} morphisms \(Y\to V\).
If \(X\to Y\) is an epimorphism in the category of presheaves,
then \(i_1=i_2\). In this case, for each \(S\), the map
\(X(S)\to Y(S)\) is surjective; since it is also injective,
it is a bijection, and thus \(X\to Y\) is an isomorphism.

Finally let us work in the category \(\widetilde{\mathcal C}\)
of sheaves. By 4.4.1 (ii), the amalgamated sum in
\(\widetilde{\mathcal C}\) of the sheaves \(Y,Y\) under \(X\)
is the sheaf \(Z\) associated to the presheaf \(V\).
Consider the diagram of morphisms:
\[
\xymatrix{
X\ar[r]^f & Y\ar@<.5ex>[r]^{i_1}\ar@<-.5ex>[r]_{i_2}
 & V\ar[r]^\tau & Z=a(V). }
\]
One has \(i_1f=i_2f\), hence \(\tau i_1f=\tau i_2f\), and thus
\(\tau i_1=\tau i_2\), since \(f\) is an epimorphism in
\(\widetilde{\mathcal C}\). By point (iii) of the lemma below,
the presheaf \(V\) is separated, hence \(\tau\) is a monomorphism
(4.3.11 (iv)). Thus \(i_1=i_2\), and one has seen above that
this implies that \(f:X\to Y\) is an isomorphism.
This proves Lemma 4.4.4, once one has verified that \(V\) is separated.
\end{proof}
